\documentstyle[%


righttag%


,11pt%


,amssymb%


,subeqn%


,russian%               remove in Eng.


,izvru%                 Set izven% in Eng.


]{amsart}





% 





\newcommand{\tts}[1]{}


% \documentstyle[a4,12pt,russian1,subeqn,mart1]{article}


\newcommand{\beq}{\begin{eqnarray}}


\newcommand{\beqv}{\begin{eqnarray*}}


\newcommand{\eeq}{\end{eqnarray}}


\newcommand{\eeqv}{\end{eqnarray*}}


\newcommand{\vf}{\varphi}


\newcommand{\La}{\Lambda}


\newcommand{\si}{\sigma}


\newcommand{\no}{\notag}


\newcommand{\pa }{\partial}


\newcommand{\eps}{\varepsilon}


\newcommand{\ovl}{\overline}


\newcommand{\Oh}[1] {{\cal O} (#1)}


\newcommand{\G}{\varGamma}


\newcommand{\om} {\omega}


\newcommand{\Db }{\ovl{D} }


\newcommand{\Dbh}{\ovl{D}_h }


\newcommand{\de}{ \delta}


\newcommand{\dxx}{\delta_{\ovl{xs}\,\widehat{xs}}}


\newcommand{\eq }{\equiv}


\newcommand{\be}{\begin{equation}{}}


\newcommand{\ee}{\end{equation}{}}


\newcommand{\beqs}{\begin{subequations}}


\newcommand{\eeqs}{\end{subequations}}


\numberwithin{equation}{section}


\def\alph{\asbuk}


%%%%%%% ^^^^^^ remove in eng.





\theoremstyle{plain}


\theorembodyfont{\it}


\newtheorem{thms}{\hskip\parindent Теорема}[section]





\theoremstyle{definition}


\theorembodyfont{\rm}


\newtheorem{notenonum}{\hskip\parindent Замечание}


\renewcommand{\thenotenonum}{}





%\hoffset=-15mm%                  For Eng.


\setcounter{page}{74}


\god{2003}


\nomer{1~(488)} %                        Remove in Eng.


%\nomer{Vol.\,47, No.\,1}%               Set in Eng.


%\str{\pageref{begin}--\pageref{end}}%  Set in Eng.


\udk{519.632}





\author{\it Г.И.\,Шишкин}


% NB:   ^^ remove in Eng.


\title{Численные методы на адаптивных сетках для сингулярно возмущенных


эллиптических уравнений\\ в области с криволинейной границей}


\thanks{Работа выполнена при финансовой поддержке Российского фонда


фундаментальных исследований (код проекта 01-01-01022).}





\date{}


%\pagestyle{myheadings}%         Set in Eng.


\pagestyle{plain} %              Remove in Eng.


\begin{document}


%\thispagestyle{plain}%          Set in Eng.


\maketitle


\label{begin}








В области $\Db$ с криволинейной границей рассматривается


краевая задача для сингулярно возмущенных эллиптических


уравнений реакции-диффузии. Классические разностные схемы


для такой задачи сходятся лишь при условии


$\eps\gg N_1^{-1}+N_2^{-1}$, где $\eps$ --- возмущающий параметр,


величины $N_1$ и $N_2$ определяют число узлов сетки по осям


$x_1$ и $x_2$. Исследуются схемы на прямоугольных сетках,


локально сгущающихся


в окрестности пограничного слоя в том случае, когда ориентация


шаблона разностной схемы в области сгущения не зависит от направления


нормали к границе. Показано, что в классе классических


разностных аппроксимаций задачи на адаптивных ``кусочно-равномерных''


сетках (сгущающихся в слое) не существует схем, сходящихся


$\eps$-равномерно и даже при условии $\eps \approx P^{-1}$,


где $P$ --- число узлов сеточной области. Таким образом, для


построения схем, сходящихся на $\Db$ при условии $\eps \le P^{-1}$,


использование сеток, подстраивающихся в погранслое к


криволинейной границе, является необходимым.


Строятся схемы, сходящиеся на $\Db$ при более слабом условии, накладываемом на


$\eps$, чем в случае традиционных разностных схем,


а также $\eps$-равномерно сходящиеся схемы, однако сходящиеся вне


пограничного слоя и на достаточно узких множествах, ``пересекающих''


пограничный слой.








\section{Введение}\label{section.1}





При исследовании процессов тепло-массообмена в средах


с малыми коэффициентами теплопроводности/диффузии возникают


краевые задачи для сингулярно возмущенных


уравнений в частных производных


(с возмущающим параметром $\eps$ --- коэффициентом $\eps^2$


при старших производных уравнений). Решения таких задач при малых


значениях параметра $\eps$ имеют особенности типа


пограничных слоев.





Ошибки решений хорошо известных численных методов существенно


зависят от величины параметра $\eps$ и могут быть большими при


малых значениях $\eps$. Таким образом, актуальной является


проблема разработки специальных численных методов для сингулярно


возмущенных задач в областях произвольной формы.


При разработке таких методов широко используются сгущающиеся


в погранслоях сетки, адаптирующиеся к границе области; шаг таких


``анизотропных'' сеток в направлении нормали к границе области зависит


от величины параметра $\eps$ и для фиксированного числа узлов


сеточной области стремится к нулю при $\eps\to 0$ (см., напр.,


\cite{1}--\cite{4} и библиографию там же).





Весьма привлекательными являются сеточные методы на локально сгущающихся


в окрестности погранслоя (априорно, либо апостериорно)


``изотропных'' сетках --- сетках, близких в области сгущения к


прямоугольным, имеющих соизмеримый шаг по всем направлениям и сохраняющих


форму шаблона разностной схемы вне зависимости от поведения границы.


Такие сетки часто применяются при уточнении сеточных решений в численных


методах на основе апостериорно адаптирующихся сеток.





В данной работе рассматриваются разностные аппроксимации для


сингулярно возмущенных эллиптических уравнений реакции-диффузии


в области с криволинейной границей. При построении схем


используются локально сгущающиеся сетки на шаблонах, ориентация которых


не связана с нормалью к границе. Строятся и исследуются


разностные схемы, сходящиеся при условии, накладываемом на параметр


$\eps$, более слабом, чем в случае традиционных разностных схем.


Обсуждаются проблемы, связанные с построением $\eps$-равномерно


сходящихся схем на указанных сетках.








\section{Постановка задачи. Обоснование исследований}


\label{section.2}





1. В области


\begin{equation}


\label{2.1}


\Db \subset \Bbb R^2


\end{equation}


с криволинейной границей $\G=\Db\setminus D$ рассмотрим


краевую задачу для сингулярно возмущенного эллиптического уравнения


реакции-диффузии


\begin{subequations}


\label{2.2}


\begin{equation}


\label{2.2a}


Lu(x) \eq\displaystyle \bigg\{\eps^2 \sum_{s=1,2} a_s(x)


\frac{\pa^2}{\pa x_s^2} - c(x)\bigg\} u(x) = f(x), \quad x\in D, 


\end{equation}


\begin{equation}


\label{2.2b}


u(x)=\vf(x), \quad x\in \G.


\end{equation}


\end{subequations}


Коэффициенты и правая часть уравнения, а также граничная функция


удовлетворяют условиям\,\footnote{\,Здесь и ниже через


$M$, $M_i$ (через $m$) обозначаются


достаточно большие (соответственно малые) положительные постоянные,


не зависящие от $\eps$ и от параметров разностных схем.


Запись $L_{(j.k)}$ (соответственно $M_{(j.k)}$ $D_{h(j.k)}$) означает, что


эти операторы (соответственно постоянные, сетки) введены в формуле $(j.k)$.}


\begin{gather*}


0<a_0\le a_s(x)\le a^0, \quad 0<c_0\le c(x)\le c^0, \quad 


x\in \Db, \quad s=1,2; \\


|f(x)|\le M, \quad x\in \Db; \quad


|\vf(x)|\le M, \quad x\in \G;


\end{gather*}


параметр $\eps$ принимает произвольные значения


из полуинтервала $(0,1]$. Граница области и функции


$a_s(x)$, $c(x)$, $f(x)$, $x\in \Db$, $\vf(x)$, $x\in \G$,


предполагаются достаточно гладкими.


При стремлении параметра $\eps$ к нулю в окрестности границы


$\G$ появляется пограничный слой.





2. Ошибки решений классических разностных схем при использовании


в качестве базовых прямоугольных сеток зависят от


величины параметра $\eps$ и становятся малыми лишь при значениях


параметра $\eps$, существенно превышающих ``эффективный'' шаг сетки.


Так, согласно оценкам (\ref{3.3}), (\ref{3.5})


классические разностные схемы (\ref{3.2}), (\ref{3.1}) и


(\ref{3.2}), (\ref{3.4}) (см. \S\,4 ниже) сходятся при условии


\begin{equation} \label{2.3}


\eps\gg N_1^{-1}+N_2^{-1},


\end{equation}


где величины $N_1$ и $N_2$ определяют число узлов сеток по переменным


$x_1$ и $x_2$ соответственно. При нарушении этого условия сеточные решения


не сходятся к решению краевой задачи (\ref{2.2}), (\ref{2.1}).





Для задачи (\ref{2.2}), (\ref{2.1}) с использованием классических


аппроксимаций краевой задачи и локально сгущающихся прямоугольных


сеток требуется построить разностные схемы, сходящиеся


$\eps$-равномерно, либо схемы, близкие к таковым. В частности,


представляют интерес схемы, сходящиеся при условии более слабом,


чем условие (\ref{2.3}).





{\em О содержании работы.}


Классическая разностная схема при использовании прямоугольных сеток,


являющихся прямым произведением сеток по $x_1$ и $x_2$,


рассматривается в \S\,4. В \S\,5 обсуждаются проблемы, связанные с


построением $\eps$-равномерно сходящихся схем на локально сгущающихся


прямоугольных сетках --- адаптивных сетках, не согласованных


с границей. Для таких сеток в \S\,6 строятся схемы, сходящиеся на


всей области $\Db$ при условии более слабом, чем условие (\ref{2.3}),


а также $\eps$-равномерно сходящиеся схемы, однако сходящиеся вне


пограничного слоя, а также на достаточно узких множествах,


пересекающих пограничный слой. В \S\,7 рассмотрена $\eps$-равномерно


сходящаяся (на всей области $\Db$) схема на сетках,


подстраивающихся к границе в окрестности пограничного слоя.


Априорные оценки, используемые в построениях, приведены в


следующем разделе.





\section{Априорные оценки}\label{section.6}





Приведем априорные оценки решения краевой задачи


(\ref{2.2}), (\ref{2.1}), используемые при построениях


(см., напр., \cite{1}, а также \cite{2}--\cite{4}).





Решение задачи представим в виде суммы функций


\begin{equation} \label{6.1}


u(x)=U(x)+V(x),\quad x\in\Db,


\end{equation}


где $U(x)$ и $V(x)$ --- регулярная и сингулярная части решения.





Для функций $u(x),\ U(x),\ V(x)$ получаются оценки


\begin{align}


\label{6.2}


&\bigg|\frac{\partial^{k}}{\partial x_1^{k_1}x_2^{k_2}}\,


u(x)\bigg| \leq M \eps^{-k},\\


%


&\bigg|\frac{\partial^{k}}{\partial x_1^{k_1}x_2^{k_2}}\,


U(x)\bigg| \leq  M, \notag \\


%


&|V(x)| \leq M\exp(-m_1\eps^{-1}r),


\quad x\in\Db, \quad k\leq 4, \notag


\end{align}


где $r=r(x,\G)$ --- расстояние от точки $x$ до множества $\G$, 


$m_1$ ---  произвольная постоянная из интервала $(0,m_0)$, 


$m_0=\min\limits_{\Db}[a^{-1}(x)c(x)]^{1/2}$.





\medskip


В $\rho$-окрестности границы $\G$, где $\rho\le m$,


выполняется оценка


\begin{gather}


\label{6.3}


\bigg|\frac{\partial^{k}}{\partial \xi_1^{k_1}\xi_2^{k_2}}\,


\widetilde{V}(\xi)\bigg| \leq M\eps^{-k_1}


\exp\big(-m_1\eps^{-1}r(\xi,\widetilde{\G})\big),\\


%


\xi\in \widetilde{\Db}, \quad r(\xi,\widetilde{\G})\le m,


\quad k\le 4, \quad m_1=m_{1(\ref{6.2})}. \notag


\end{gather}


Здесь $\xi=\xi(x)$ --- ортогональная система координат, связанная с


границей $\G$; $\xi_1$ --- переменная вдоль нормали к границе $\G$;


$\widetilde{V}(\xi)=V(x(\xi))$; $\widetilde{D}$ и


$\widetilde{\G}$  --- образы множеств $D$ и $\G$; 


$\xi_1=r(\xi,\widetilde{\G})$.





\begin{thms} \label{t.6.1}


Пусть $a_s, c, f \in C^{l+\alpha}(\Db)$, $\vf \in C^{l-1+\alpha}(\G)$, 


$\G\in C^{l-1+\alpha}$, $s=1,2$, $l=6$, $\alpha>0$. Тогда для решения


краевой задачи $(\ref{2.2})$, $(\ref{2.1})$ и его компонент из


представления $(\ref{6.1})$ справедливы оценки $(\ref{6.2})$,


$(\ref{6.3})$.


\end{thms}





\section{Классическая разностная схема}\label{section.3}





Выпишем классическую разностную схему для задачи (\ref{2.2}),


(\ref{2.1}) и обсудим некоторые трудности, возникающие при


численном решении задачи, когда параметр $\eps$ достаточно мал.





1. Рассмотрим простейшую монотонную разностную схему из


(\cite{5}, c.\,218). Пусть на $\Bbb R^2$ введена (базовая) сетка


\beqs \label{3.1}


\begin{equation}\label{3.1a}


\Db{}_h^{(2)}=\om_1 \times \om_2,


\end{equation}


где $\om_s$ --- сетка на оси $x_s$, вообще говоря, неравномерная,


$s=1,2$; считаем выполненным условие $h\le MN^{-1}$.


Здесь $h=\max\limits_s h_s$, $h_s=\max\limits_i h_s^i$,


$h^i_s=x_s^{i+1}-x_s^i$, $x_s^i, x_s^{i+1} \in \om_s$,


$N=\min\limits_s N_s$, $N_s+1$ --- минимальное число узлов


сетки $\om_s$ на отрезке единичной длины. На $\Db$ введем сетку


\begin{equation} \label{3.1b}


\Dbh=D_h \cup \G_h, \quad \Dbh=\Dbh\big(D^{(2)}_{h(\ref{3.1a})}\big)=


\Dbh\big(D^{(2)}_{h(\ref{3.1a})};\, D_{(\ref{2.1})}\big),


\end{equation}


\eeqs


где $D_h=D\cap D^{(2)}_h$, множество $\G_h$ образовано пересечением


прямых $x_s=x_s^i$, $x_{3-s}\in \Bbb R$, $s=1,2$, проходящих через


узлы сетки $D^{(2)}_h$, с границей $\G$.





Задачу (\ref{2.2}), (\ref{2.1})


аппроксимируем разностной схемой (\cite{5}, c.\,221)


\begin{equation}\label{3.2}


\begin{split}


\La z(x) &\eq \bigg\{ \eps^2 \sum_{s=1,2} a_s(x) \dxx -c(x)\bigg\}


z(x) = f(x), \quad x\in D_h, \\


z(x)&= \vf(x), \quad x \in \G_h.


\end{split}


\end{equation}


Здесь


$\dxx\,z(x)$ --- вторая (центральная) разностная производная на


неравномерной сетке, например, 


$\de_{\ovl{x1}\,\widehat{x1}}\,z(x)=2 (h_1^i+h_1^{i-1})^{-1}


[\delta_{x1}\,z(x)-\delta_{\ovl{x1}}\,z(x)]$,


$x=(x_1^i,x_2) \in D_h$, 


$\delta_{x1}\,z(x)$ и $\delta_{\ovl{x1}}\,z(x)$ --- первые разностные


производные (вперед и назад).





Для разностной схемы (\ref{3.2}), (\ref{3.1}) справедлив


принцип максимума (\cite{5}, c.\,228).





Принимая во внимание априорные оценки решения краевой задачи,


находим оценку


\begin{equation} \label{3.3}


|u(x)-z(x)|\le M(\eps+N^{-1})^{-1}N^{-1},\quad 


x\in \Db_{h(\ref{3.1})}.


\end{equation}


В случае сетки


\begin{equation} \label{3.4}


\Dbh=\Dbh(D_h^{(2)}),


\end{equation}


где $D_h^{(2)}$ --- равномерная прямоугольная сетка,


имеем оценку


\begin{equation} \label{3.5}


|u(x)-z(x)|\le M(\eps+N^{-1})^{-2}N^{-2},\quad 


x\in \Db_{h(\ref{3.4})}.


\end{equation}


Из рассмотрения модельных задач вытекает, что схемы (\ref{3.2}), (\ref{3.1})


и (\ref{3.2}), (\ref{3.4}) не сходятся $\eps$-равномерно.





2. Обсудим условия, при которых решение разностной схемы сходится


при $N\to\infty$ и $\eps\to 0$, где


$\eps\to 0$ при $N\to\infty$.





Введем ряд определений. Пусть $z(x)$, $x\in\Dbh$, ---


решение некоторой разностной схемы. Скажем, что оценка


$$


|u(x)-z(x)|\le M[(\eps+N_1^{-1})^{-\nu_1^1}


N_1^{-\nu_1}+(\eps+N_2^{-1})^{-\nu_2^1}N_1^{-\nu_2}],


\quad x\in\Dbh,


$$


где  $\nu_i, \nu^1_i \ge 0$, неулучшаема по вхождению величин


$N_1$, $N_2$, $\eps$, если оценка


$$


|u(x)-z(x)|\le M[(\eps+N_1^{-1})^{-\alpha_1^1}


N_1^{-\alpha_1}+(\eps+N_2^{-1})^{-\alpha_2^1}N_1^{-\alpha_2}


], \quad x\in\Dbh,


$$


вообще говоря, неверна при выполнении условия


$\alpha_i\ge \nu_i$, $\alpha_i^1\le \nu_i^1$, причем


$\alpha_1+\alpha_2-\alpha_1^1-\alpha_2^1>


\nu_1+\nu_2-\nu_1^1-\nu_2^1$.


В случае, когда найдется функция


$\mu\big(N_1^{-1},\,N_2^{-1}\big)$, сходящаяся к нулю


$\eps$-равномерно при $\,N\to\infty$, такая, что


выполняется оценка


$$


|u(x)-z(x)|\le M\min[


\mu(\eps^{-\nu} N_1^{-1}, \, \eps^{-\nu} N_2^{-1}),\,


1], \quad x\in\Dbh,


$$


где $\nu>0$ --- некоторая постоянная, будем говорить, что схема


сходится с дефектом $\Oh{\eps^{-\nu}}$. Если же величина $\nu$


может быть сколь угодно малой, будем говорить, что схема сходится


почти $\eps$-равномерно (с дефектом сходимости $\Oh{\eps^{-\nu}}$).


При $\nu=0$ схема сходится $\eps$-равномерно.





Оценки (\ref{3.3}), (\ref{3.5}) неулучшаемы по вхождению


величин $N$, $\eps$.  В силу неулучшаемости оценок


(\ref{3.3}) и (\ref{3.5}) порядок $\nu$ дефекта схем


(\ref{3.2}), (\ref{3.1}) и (\ref{3.2}), (\ref{3.4}) неулучшаем


и равен единице. Эти схемы сходятся при условии


$N_1^{-1}, N_2^{-1} \ll \eps$:


\begin{equation} \label{3.6}


\eps^{-1}=o(N_1), \quad \eps^{-1}=o(N_2).


\end{equation}


Условие (\ref{3.6}) является необходимым и достаточным для сходимости схем


(\ref{3.2}), (\ref{3.1}) и (\ref{3.2}), (\ref{3.4}).





\begin{thms} \label{t.3.1}


Пусть для решения краевой задачи $(\ref{2.2})$, $(\ref{2.1})$


и его компонент из представления $(\ref{6.1})$ выполняются


априорные оценки $(\ref{6.2})$, $(\ref{6.3})$.


В случае разностных схем {\em (\ref{3.2}), (\ref{3.1})} и


{\em (\ref{3.2}), (\ref{3.4})} условие $(\ref{3.6})$


является необходимым и достаточным для сходимости


сеточных решений к решению краевой задачи $(\ref{2.2})$,


$(\ref{2.1})$ при $N\to\infty$ и $\eps\to 0;$


дефект сходимости схем есть $\Oh{\eps^{-1}}$.


Для сеточных решений выполняются оценки $(\ref{3.3})$, $(\ref{3.5});$


оценки неулучшаемы по вхождению величин $N$, $\eps$.


\end{thms}








\section{О построении $\eps$-равномерно сходящихся


схем\protect\\ на локально сгущающихся сетках}\label{section.4}





Заметим, что особенность типа пограничного слоя


при удалении от границы $\G$ экспоненциально


убывает с показателем экспоненты $m\eps^{-1} r(x,\G)$ (см.


последнюю оценку в (\ref{6.2}) и оценку (\ref{6.3})).


Сингулярная компонента $V(x)$ из (\ref{6.1}) при


$r(x,\G)\ge \si$ не превосходит величины $M\de$, где


$\de$ --- достаточно малое число, когда


$\si=m_1^{-1}\eps\ln \de^{-1}$,


$m_1=m_{1(\ref{6.2})}$.


Невязка для разностной схемы на решении краевой задачи велика


(однако $\eps$-равномерно конечная) лишь в окрестности погранслоя,


являющейся достаточно узкой при малых значениях параметра $\eps$.





Обсудим проблемы, возникающие при построении $\eps$-равномерно


сходящихся схем.





1. Имея в виду использование схем на локально сгущающихся


сетках при решении краевой задачи, введем некоторые


характеристики схем (сеток).


Через $P$ ({\it объем} сеточной задачи) обозначим число узлов


сетки на $\Db$, в которых требуется найти сеточное решение.


В случае локально сгущающихся сеток будет удобно оценивать


сходимость сеточных решений через величину $P^{-1/n}$, где


$n=2$ --- размерность пространства.





В случае сеток (\ref{3.1}) для величины $P$ выполняется оценка


$$


m N_1N_2 \le P\le M N_1N_2.


$$


Для решений схемы (\ref{3.2}) на сетках (\ref{3.1}) и (\ref{3.4})


имеем оценки


\begin{alignat*}{2}


|u(x)-z(x)|&\le M\min[\eps^{-1}(N_1^{-1}+


N_1P^{-1}), 1], &\quad x&\in\Db_{h(\ref{3.1})};\\


%


|u(x)-z(x)|&\le M\min[\eps^{-2}(N_1^{-1}+


N_1P^{-1})^2, 1], &\quad x&\in\Db_{h(\ref{3.4})}.


\end{alignat*}





Будем говорить, что распределение узлов сетки по $x_1$ и $x_2$


сбалансировано (или короче --- сетка сбалансирована), если


на ней достигается оптимальный при фиксированном значении параметра


$\eps$ порядок скорости сходимости. Схему на сбалансированной сетке


будем также называть сбалансированной.





Для решений схемы (\ref{3.2}) на сбалансированных сетках


(\ref{3.1}) и (\ref{3.4}) получаются оценки


\begin{alignat}{2}


\label{4.1}


|u(x)-z(x)|&\le M(\eps+P^{-1/2})^{-1}


P^{-1/2}, &\quad x&\in\Db{}^{b}_{h(\ref{3.1})}; \\


%


\label{4.2}


|u(x)-z(x)|&\le M(\eps^2+P^{-1})^{-1}


P^{-1}, &\quad x&\in\Db{}^{b}_{h(\ref{3.4})},


\end{alignat}


где $\Db{}_h^{b}$ --- сбалансированная сетка. Дефект сходимости схем


(\ref{3.2}), (\ref{3.1}) и (\ref{3.2}), (\ref{3.4})


относительно величины $P^{1/2}$ на сбалансированных сетках


есть величина $\Oh{\eps^{-1}}$ (для


несбалансированных сеток дефект больше). Эти схемы сходятся при


условии  $P^{-1/2} \ll \eps$:


\begin{equation} \label{4.3}


\eps^{-1}=o(P^{1/2}).


\end{equation}


Это условие является необходимым и достаточным для сходимости схем.





2. Рассмотрим класс разностных схем, основанных на классических


аппроксимациях краевой задачи в случае ``кусочно-равномерных'' сеток,


а именно, на локально сгущающихся сетках, являющихся равномерными


в ближайшей окрестности границы $\G$ и вне некоторой несколько


большей окрестности. Сетки, вообще говоря, не согласованы с границей


области.





Для простоты пусть $x_1+x_2=0$ есть уравнение границы $\G$ в единичном


круге (с центром в точке $(0,0)$) и область $\Db$ содержит часть круга,


для которой $x_1+x_2 \ge 0$.


Пусть каким-то образом построена сетка


\begin{equation} \label{4.4}


\Db{}_h^{*}=\Db{}_h^{*}(\rho_1),


\end{equation}


где $\rho_1>0$ --- параметр, выбираемый ниже, который определяет


распределение узлов сетки (\ref{4.4}). Эта сетка


равномерна на каждом из множеств


$D_1=D_1(\rho_1)\,$ и


$\,D_2=D\setminus \Db_1(M\rho_1)$, где 


$$


D_1(\rho_1)=\{x: r(x,\G)<\rho_1\}


$$


--- $\rho_1$-окрестность множества $\G$; $\rho_1\le 4^{-1}$.


Сетки $D_{ih}=D_i\bigcap \Db{}^{*}_{h(\ref{4.4})}$


имеют шаг $h^i_s$ по оси $x_s$; $i,s=1,2$.


Для простоты считаем, что шаблоны пятиточечных


схем с центром шаблона из $D_{ih}$ являются правильными,


т.е. их левые и правые, а также верхние и нижние плечи одинаковы


и равны $h_1^{(i)}$ и $h_2^{(i)}$ соответственно, $i=1,2$.





Рассмотрим фрагменты сеточной задачи из класса разностных схем


на сетках (\ref{4.4}), а именно, фрагменты на множествах


$\Db_{1h}$ и $\Db_{2h}$.


Пусть $z_i(x)$, $x\in \Db_{ih}$, --- решение задачи


\begin{equation} \label{4.5}


\begin{split}


\La_{(\ref{3.2})} z_i(x) &= f(x), \quad 


x\in D_{ih}, \\


z_i(x)&=u(x), \quad x\in \G_{ih}, \qquad


i=1,2,


\end{split}


\end{equation}


где $u(x)$, $x\in\Db$, --- решение задачи (\ref{2.2}), (\ref{2.1}).





Для функций $z_i(x)$, $x\in \Db_{ih}$, имеем оценки


\beqs \label{4.6}


\beq


% && \no \\[-6ex]


&& \label{4.6a} \hspace*{-1cm}


|u(x)-z_1(x)|\le M[


(\eps+h_1^{(1)})^{-2} (h_1^{(1)})^2+


(\eps+h_2^{(1)})^{-2} (h_2^{(1)})^2], \quad x\in \Db_{1h}; \\[1.5ex]


%


&& \hspace*{-1cm} \no


|u(x)-z_2(x)|\le M\{[


(\eps+h_1^{(2)})^{-2} (h_1^{(2)})^2+


(\eps+h_2^{(2)})^{-2} (h_2^{(2)})^2]\times \\[1ex]


&& \label{4.6b}


\hspace*{3cm} 


\times \max\limits_{\G_{2h}}|V(x)| +


(h_1^{(2)})^2 + (h_2^{(2)})^2\}, \quad


x\in \Db_{2h}, \hspace*{1cm}


\eeq


\eeqs


где $V(x)=V_{(\ref{6.1})}(x)$ --- сингулярная компонента решения;


оценки (\ref{4.6a}) и (\ref{4.6b}) неулучшаемы по вхождению величин


$h_1^{(1)}$, $h_2^{(1)}$, $\eps$ и


$h_1^{(2)}$, $h_2^{(2)}$, $\eps$ соответственно.





Анализ ошибок $u(x)-z_i(x)$, $x\in \Db_{ih}$, при $\rho_1\approx \eps$


показывает, что для сбалансированных схем не существует


кусочно-равномерных сеток $\Db{}^{*}_{h(\ref{4.4})}$, на которых


решения задач (\ref{4.5}) сходятся $\eps$-равномерно к решению задачи


(\ref{2.2}), (\ref{2.1}). В случае задач (\ref{4.5}), (\ref{4.4})


не существует сеток, на которых дефект сходимости есть


$\Oh{\eps^{-1/2}}$.


Подобное поведение ошибок для сеточных решений сохраняется


и в случае семейства сеток


\begin{equation} \label{4.7}


\Db{}_h^{*},


\end{equation}


являющихся прямоугольными и равномерными в $\rho$-окрестности


конечного участка границы $\G$, где $\rho=\rho(\eps,P)$, причем


$\eps^{-1}\rho\to \infty$ при $P\to\infty$.





\begin{thms} \label{t.4.1}


Для краевой задачи $(\ref{2.2})$, $(\ref{2.1})$


в классе сбалансированных разностных схем, строящихся на


основе классических аппроксимаций краевой задачи на


локально сгущающихся сетках вида $(\ref{4.7})$ $($не согласованных


с границей $\G)$, не существует $\eps$-равномерно сходящихся схем$;$


более того, не существует схем, дефект сходимости которых есть


$\Oh{\eps^{-1/2}}$.


\end{thms}





\begin{note}


Условие $P^{-1}\ll \eps$:


\begin{equation} \label{4.8}


\eps^{-1}=o(P),


\end{equation}


вообще говоря, не достаточно для сходимости схемы, если сетка не


согласована с границей области.


\end{note}





\begin{note}


Как следует из приведенных построений,


использование локально сгущающихся сеток не позволяет существенно


ослабить условие (\ref{3.6}) сходимости классических разностных


схем. Применением достаточно общих локально сгущающихся сеток,


однако не согласованных с границей, невозможно уменьшить уже в два раза


порядок $\nu=1$ дефекта сходимости классических схем.


\end{note}





\begin{note}


Для построения $\eps$-равномерно сходящихся на


$\Db$ разностных схем на основе классических аппроксимаций краевой


задачи использование сеток, ориентация шаблонов которых в


погранслое согласована с направлением нормали к границе области,


является необходимым. Более того, такие согласованные сетки


необходимо использовать и для построения схем, сходящихся при условии


не слабее, чем (\ref{4.8}).


\end{note}


\addtocounter{note}{-3}








\section{Разностные схемы на сетках, не согласованных с


границей}\label{section.5}





Как следует из результатов \S\,5, разностные схемы на


прямоугольных сетках, в том числе локально сгущающихся,


в случае, когда сетки не согласованы с границей, не позволяют


получать решения, сходящиеся $\eps$-равномерно на всей сеточной


области $\Dbh$. В этой связи возникает проблема построения разностных


схем, сходящихся $\eps$-равномерно на подмножествах из $\Db$,


в частности, на ``сечениях'' множества $\Db$ --- на связных подмножествах,


проходящих через пограничный слой.





1. Рассмотрим область сходимости разностной схемы


(\ref{3.2}), (\ref{3.4}).


Пусть $v(x)$, $x\in\Db$, --- достаточно гладкая функция.


Через $v^h(x)$, $x\in\Dbh$, где $\Dbh=\Db_{h(\ref{3.1})}$,


обозначим  решение задачи


\begin{alignat*}{2}


\La_{(\ref{3.2})} z(x) &= f(x), &\quad 


x&\in D_{h},\\


z(x)&=v(x), &\quad x&\in \G_{h},


\end{alignat*}


где $f(x)=L_{(\ref{2.2})}\,v(x)$.





Для компонент $U^h(x)$, $V^h(x)$, $x\in\Dbh$, отвечающих


компонентам $U(x)$, $V(x)$, $x\in\Db$, из представления (\ref{6.1}),


имеем оценки


\begin{gather} \label{5.1}


|U(x)-U^h(x)|\le M[N_1^{-2}+N_2^{-2}], \quad 


x\in \Db_{h}; \\


%


\left.\begin{matrix}|V(x)| \\ |V^h(x)| \end{matrix}\right\}


\le M[\exp(-m_1(\eps+N_1^{-1})^{-1} r_1) +


\exp(-m_2(\eps+N_2^{-1})^{-1} r_2)], \quad


x\in \Db_{h}, \notag


\end{gather}


где $r_s=r_s(x,\G)$ --- расстояние, измеряемое вдоль оси $x_s$,


между точкой $x$ и множеством $\G$, $s=1,2$;


$m_1$, $m_2$ --- достаточно малые постоянные,


$m_1, m_2< 2^{-1/2} m_{1(\ref{6.2})}$. Из оценок (\ref{5.1})


вытекает $\eps$-равномерная сходимость решений схемы


(\ref{3.2}), (\ref{3.4}) вне $\rho_1$--окрестности множества $\G$


\begin{equation} \label{5.2}


|u(x)-z(x)|\le M[ N_1^{-2}+N_2^{-2}], \quad


x\in \Db_{h}, \ \ r(x,\G)\ge \rho_1,


\end{equation}


где


\begin{gather} \label{5.3}


\rho_1=\rho_1(\eps,N_1,N_2)=\min\{ M_1[


(\eps+N_1^{-1})\ln N_1+


(\eps+N_2^{-1})\ln N_2], m\}, \\


M_1=2(m_1^{-1}+m_2^{-1}),\ \


m_i=m_{i(\ref{5.1})}. \notag


\end{gather}


При условии $\eps\ll \ln^{-1} N_1+\ln^{-1} N_2$, 


$N_1, N_2\to\infty\,$ величина $\rho_1$ стягивается к нулю.





2. Построим разностную схему, сходящуюся при условии более слабом, чем


условие (\ref{4.3}).


Через $z_0(x)$, $x\in\Dbh,$ обозначим решение задачи


(\ref{3.2}), (\ref{3.4}). Пусть


\beqs \label{5.4}


\begin{equation} \label{5.4a}


D_1 


\end{equation}


--- $\rho_1$-окрестность множества $\G$; $\rho_1=\rho_{1(\ref{5.3})}$.


На множестве $\Db_1$ строим равномерную прямоугольную сетку


\begin{equation} \label{5.4b}


\Db_{1h}=\Db_{(\ref{3.1b})}(D_h^{(2)};\, D_1),


\end{equation}


где $D_{h(\ref{5.4b})}^{(2)}$ --- равномерная сетка с шагом $h_s^{(2)}$


по оси $x_s$, определяемым соотношением 


$h_1^{(2)}=h_2^{(2)}=\rho_1^{1/2} N_1^{-1/2} N_2^{-1/2}$. 


Заметим, что для числа узлов сетки $\Db_{1h}$ выполняется оценка


$P(\Db_{1h})\le MN_1N_2$.





На сетке $\Db_{1h}$ решаем задачу


%\beq \label{5.4c}


%&& \La_{(\ref{3.2})} z_1(x) = f(x), \quad 


%x\in D_{1h}, \hspace*{1cm} \no\\ [-1.5ex]


%&& \\[-1ex]


%&& \phantom{\La_{(\ref{3.2})}\,}


%z_1(x)= \left\{


%\begin{array}{ll}


%\vf(x), & \ x\in \G_{1h}\bigcap \G,\\[1ex]


%\ovl{z}_0(x), & \ x\in \G_{1h}\setminus \G,


%\end{array} \right. \no


%\eeq


%


\begin{equation}\label{5.4c}


\begin{split}


\La_{(\ref{3.2})} z_1(x) &= f(x), \quad 


x\in D_{1h}, \\


z_1(x)&=


\begin{cases}


\vf(x), & x\in \G_{1h}\bigcap \G;\\


\ovl{z}_0(x), & x\in \G_{1h}\setminus \G,


\end{cases}


\end{split}


\end{equation}


где $\ovl{z}_0(x)$, $x\in\Db$, --- некоторый интерполянт,


строящийся по значениям $z_0(x)$, $x\in\Dbh$.


При построении $\ovl{z}_0(x)$ минимальный многоугольник


$\ovl{D}{}_h^{(m)}$, содержащий узлы сетки


$\Dbh$, разбивается прямыми


$x_1=x_1^i$ и $x_2=x_2^j$, $(x_1^i, x_2^j)\in


\ovl{D}{}_{h(\ref{3.4})}^{(2)}$ на ячейки.


На прямоугольных ячейках функция $\ovl{z}_0(x)$ есть билинейный


интерполянт. Непрямоугольные ячейки триангулируются.


Функция $\ovl{z}_0(x)$ на таких треугольных ячейках является


линейной. Так построенные билинейные и линейные функции


$\ovl{z}_0(x)$ продолжаются на ближайшие ``криволинейные'' элементы


из множества $\Db\setminus \ovl{D}{}_h^{(m)}$.


Функция $\ovl{z}_0(x)$ непрерывна на $\Db$.





Функцию


\begin{equation} \label{5.4e}


z^*(x)= \left\{


\begin{array}{ll}


z_1(x), & x\in \Db_{1h},\\[2mm]


z_0(x), & x\in \Db_{h(\ref{3.4})}\setminus \Db_1 


\end{array} \right\}, \quad \ x\in \Db{}_h^{*},


\end{equation}


где


\begin{equation} \label{5.4d}


\Db{}_h^{*}=\Db_{1h} \bigcup  \{\Db_{h(\ref{3.4})}


\setminus \Db_1 \},


\end{equation}


\eeqs


назовем решением разностной схемы


(\ref{3.2}), (\ref{3.4}), (\ref{5.4}).





Для решения этой схемы получается оценка


\begin{multline} \label{5.5}


|u(x)-z^*(x)|\le M\big(\eps^2 + N_1^{-1}N_2^{-1}


(N_1^{-1}\ln N_1 + N_2^{-1}\ln N_2)\big)^{-1}\times \\


%


\times \big[\eps


N_1^{-1}N_2^{-1}(\ln N_1 +\ln N_2)+N_1^{-1}N_2^{-1}


(N_1^{-1}\ln N_1 + N_2^{-1}\ln N_2)\big],


\ \ x\in\Dbh^{\,*}.


\end{multline}


На сбалансированной сетке $\Dbh^{\,*b}$ имеем оценку


\begin{equation} \label{5.6}


|u(x)-z^*(x)|\le M(\eps^{-2}+


P^{-3/2}\ln P)^{-1}[


\eps P^{-1}\ln P + P^{-3/2}\ln P],


\ \ x\in \Db{}^{*b},


\end{equation}


где $P$ --- число узлов сетки $\Dbh^{\,*b}$.


Для сбалансированной схемы (\ref{3.2}), (\ref{3.4}), (\ref{5.4})


дефект сходимости схемы есть $\Oh{\eps^{-2/3}\ln^{1/3}\eps^{-1}}$.


Схема сходится при условии 


$P^{-1/2} \ll \eps^{2/3}\ln^{-1/3}\eps$:


\begin{equation} \label{5.7}


\eps^{-1}=o(P^{3/4}\ln^{-1/2}P)


\end{equation}


--- более слабом, чем условие (\ref{4.3}). Условие


(\ref{5.7}) неулучшаемо; при нарушении его схема


(\ref{3.2}), (\ref{3.4}), (\ref{5.4}) на сбалансированных сетках,


вообще говоря, не сходится.





\begin{thms} \label{t.5.1}


Пусть выполняется условие теоремы $\ref{t.3.1}$. В случае разностной схемы


{\em (\ref{3.2}), (\ref{3.4}), (\ref{5.4})} на сбалансированных сетках


условие $(\ref{5.7})$ является необходимым и достаточным для сходимости


сеточного решения к решению краевой задачи $(\ref{2.2})$, $(\ref{2.1});$


дефект сходимости схемы есть $\Oh{\eps^{-2/3}\ln^{1/3}\eps^{-1}}$.


Для сеточных решений выполняются оценки $(\ref{5.5})$, $(\ref{5.6});$


оценка $(\ref{5.6})$ неулучшаема по вхождению величин $P$, $\eps$.


\end{thms}





3. Построим разностную схему, сходящуюся $\eps$-равномерно на


некотором подмножестве из $\Db$, проходящем через пограничный слой.


Для простоты будем считать, что в единичном круге, имеющем центром


точку $(0,0)$, расположенную на границе $\G$, единичная (внутренняя)


нормаль к границе $n=n(x)=(n_1,n_2)$ удовлетворяет условию 


$n_1\ge m_0$. На множестве $D$ из единичного круга рассмотрим


множества


\beqs \label{5.8}


\begin{equation} \label{5.8a}


D_2 \text{ \ и \ } D_2^0


\end{equation}


--- соответственно $\rho_2$- и $2^{-1}\rho_2$-окрестности оси $x_1$; здесь


\begin{equation} \label{5.8b}


\rho_2=\rho_2(\eps,N_1,N_2)=\min[m,\,


4m_1^{-1}\eps\ln N],\quad m_1=m_{1(\ref{6.2})};


\end{equation}


$N_1$, $N_2$ определяют число узлов сетки $\Db_{h(\ref{3.1a})}$,


$N=\min[N_1, N_2]$. Пусть


\begin{equation} \label{5.8c}


D_h=\om_1\times \om_2


\end{equation}


--- кусочно-равномерная сетка на плоскости $\Bbb R^2$; шаг сетки


$\om_2$  постоянен и равен $h_2=\rho_2N_2^{-1}$, шаг сетки


$\om_1$  постоянен при $x_1< \si$ и при $x_1> \si$ и равен


$h_1^{(1)}=4\si N_1^{-1}$ и $h_1^{(2)}=2 N_1^{-1}$ соответственно;


$\si=\si(\eps,N_1,N_2)=\min[m, \rho_3+m_0(1-m_0^2)^{-1/2}\rho_2]$, 


$\rho_3=\min[m,\, 2m_1^{-1}\eps\ln N_1]$.


На множестве $\Db_{2(\ref{5.8a})}$ строим сетку


\begin{equation} \label{5.8d}


\Db_{2h}=\Db_{h(\ref{3.1b})}(D_{h(\ref{5.8c})}; \,


D_{2(\ref{5.8a})}).


\end{equation}





Краевую задачу (\ref{2.2}) на множестве $\Db_2$ аппроксимируем


разностной схемой


\begin{equation} \label{5.8e}


\begin{split}


\La_{(\ref{3.2})} z_2(x) &= f(x), \quad 


x\in D_{2h}, \\


z_2(x)&=


\begin{cases}


\vf(x), & x\in \G_{2h}\bigcap \G;\\


\ovl{z}_0(x), & x\in \G_{2h}\setminus \G,


\end{cases}


\end{split}


\end{equation}


где $z_0(x)$, $x\in\Dbh$, --- решение задачи (\ref{3.2}), (\ref{3.4}).





Далее введем функцию


\begin{equation} \label{5.8f}


z^*(x)= \left\{


\begin{array}{ll}


z_2(x), & x\in \Db_{2h},\\[1ex]


z_0(x), & x\in \Db_{h(\ref{3.4})}\setminus \Db_2 \


\end{array} \right\}, \quad x\in \Dbh^{\,*},


\end{equation}


\eeqs


где


$$


\Db{}_h^{*}=\Db_{2h}\bigcup \{\Db_{h(\ref{3.4})}\setminus \Db_2\},


$$


которую назовем решением разностной схемы (\ref{3.2}),


(\ref{3.4}), (\ref{5.8}).





Заметим, что функция (\ref{5.8f}), вообще говоря,


не сходится $\eps$-равномерно на  $\Dbh^{\,*}$, однако


сходится $\eps$-равномерно на подмножестве


$\Db{}_h^{*0}$ из  $\Db{}_h^{*}$, где


\begin{equation} \label{5.9}


\Db{}_h^{*0}=\big\{\Db^{\,0}_2 \bigcup


\{\Db_{(\ref{2.1})}\setminus D_{1(\ref{5.4a})}\}\big\}


\bigcap \Db{}_h^{\,*}.


\end{equation}





С использованием техники работы \cite{1} устанавливается оценка


\begin{equation} \label{5.10}


|u(x)-z^*(x)|\le M[N_1^{-2}\ln^2 N_1 +


N_2^{-2}\ln^2 N_2], \quad x\in \Db{}_h^{\,*0}.


\end{equation}





\begin{thms} \label{t.5.2}


Пусть выполняется условие теоремы $\ref{t.3.1}$. Тогда решение разностной


схемы {\em (\ref{3.2}), (\ref{3.4}), (\ref{5.8})} сходится на множестве


$\Db{}^{*0}_{h(\ref{5.9})}$ \ $\eps$-равномерно. Для сеточного решения


справедлива оценка $(\ref{5.10})$.


\end{thms}





\begin{notenonum}


``Суммарная'' ширина всех перемычек,


соединяющих внешность погранслоя с границей, на которых


можно достичь $\eps$-равномерную сходимость сеточных решений,


стремится к нулю при $\eps\to 0$.


\end{notenonum}





\section{Разностные схемы на шаблоне,\protect\\ подстраивающемся


в погранслое к границе области}\label{section.7}





В \cite{1} в случае областей с криволинейной границей построены схемы,


сходящиеся $\eps$-рав\-но\-мер\-но. При построении схем исходная область


разбивалась на перекрывающиеся подобласти. Подходящей заменой


координат граница


области (на подобласти разбиения) переводилась в прямую (в случае


двумерной области) и строилась прямоугольная кусочно-равномерная


сетка, сгущающаяся в погранслое. Таким образом, локально сгущающаяся


сетка подстраивалась к границе на $m$-окрестности границы $\G$.


Получающаяся сеточная задача решалась методом Шварца. Недостатком


таких разностных схем является необходимость использования итераций


при решении сеточных задач.





Представляют определенный интерес безытерационные разностные схемы,


сходящиеся $\eps$-равномерно. Приведем такую схему.


\medskip





1. Решение задачи  (\ref{3.2}), (\ref{3.4}) на множестве


$\Dbh\setminus D_{1(\ref{5.4a})}$ сходится $\eps$-равномерно


с оценкой (\ref{5.2}). Чтобы уточнить решение задачи в погранслое,


построим на $\Db_{1(\ref{5.4a})}$ согласованную с границей сетку


и соответствующую разностную схему.





В $\Db_1$ каждой точке $x$ соответствует некоторое расстояние


до границы $\G$ --- величина $r(x,\G)$  и длина дуги


(по контуру границы $\G$) от некоторой фиксированной точки до основания


нормали (к границе), проходящей через точку $x$, --- величина $l(x)$.


Перейдем в $\Db_1$ к переменным $X=X(x)$:


\begin{equation} \label{7.1}


X_1=r(x,\G), \quad X_2=l(x).


\end{equation}


В этих переменных граница $\G$ переходит в прямую. Вторая часть


границы множества $D_1$, а именно, множество $\G_1\setminus \G$


также переходит в прямую. Считаем, что множество $\Db_1$ достаточно


узкое и система (\ref{7.1}) разрешима относительно $x_1$, $x_2$.


Через $X^{-1}(x)$ обозначим отображение, обратное $X(x)$.





Для функций $v(x)$, $W(X)$ и подобластей $D^0\subseteq \Db_1$


будем использовать обозначения


\begin{gather*}


v(x(X))=v_X(X)=\{v(x)\}_X, \quad W(X(x))=W_{X^{-1}}(x), \\


D^0_X=X(D^0)=\{X: X^{-1}(x)\in D^0\}.


\end{gather*}


Полагаем


$$


\widetilde{D}_{X^{-1}\,}=X^{-1}(\widetilde{D})=\{x: X(x)\in


\widetilde{D}\},


$$


где $\widetilde{D}$ --- некоторое подмножество из $\Bbb R^2$ такое, что


$\widetilde{D}_{X^{-1}}\,\subseteq \Db_1$. Множество $\Db_{1X}$


есть вертикальная полоса.





Задача (\ref{2.2}), (\ref{5.4a}), имеющая вид


\begin{equation} \label{7.2}


\begin{split}


L u(x) &= f(x), \quad 


x\in D_{1}, \\


u(x)&=


\begin{cases}


\vf(x), & x\in \G_{1}\bigcap \G;\\[1ex]


u^*(x), & x\in \G_{1}\setminus \G,


\end{cases}


\end{split}


\end{equation}


где $u^*(x)$ --- решение задачи


(\ref{2.2}), (\ref{2.1}) при $x\in \G_{1}\setminus \G$,


в новых переменных переходит в задачу


\begin{equation} \label{7.3}


\begin{split}


L_X U(X) &= F(X), \quad 


X\in D_{1X}, \\


U(X)&=


\begin{cases}


\Phi(X), & X\in \G_{1X}\bigcap \G_X; \\


U^*(X), & X\in \G_{1X}\setminus \G_X.


\end{cases}


\end{split}


\end{equation}


Здесь


$L_X \equiv \eps^2L^2_X+L_X^0$, $L_X^0=-C(X)$,


$L_X^2 \equiv \sum\limits_{s,k=1,2} A_{sk}(X) \frac{\partial^2}


{\partial X_s \,\partial X_k} +\sum\limits_{s=1,2} B_s(X)


\frac{\partial}{\partial X_s}$.


Коэффициенты $A_{sk}$, $B_s$ определяются по формулам


\begin{gather*}


A_{sk}(X)=\bigg\{a_1(x)\frac{\partial}{\partial x_1}X_s\,


\frac{\partial}{\partial x_1} X_k+


a_2(x)\frac{\partial}{\partial x_2}X_s\,


\frac{\partial}{\partial x_2} X_k\bigg\}_X, \ \ s,k=1,2;\\


%


B_{s}(X)=\bigg\{a_1(x)\frac{\partial^2}{\partial x_1^2}\,X_s+


a_2(x)\frac{\partial^2}{\partial x_2^2}\,X_s\bigg\}_X, \ \ s=1,2.


\end{gather*}


Функции $C$, $F$, $\Phi$, $U^*$, $U$ определяются соотношением


$V(X)=v_X(X)$. Задача (\ref{7.3}) является периодической на множестве


$\Db_{1X}$ \ с периодом $l_0$.





На множестве $\Db_{1X}$  строим равномерную сетку


\begin{equation} \label{7.4}


\Db_{1Xh}=\ovl{\om}_{1X1} \times \om_{1X2}\ ,


\end{equation}


где $\ovl{\om}_{1X1}\,$ и $\,\om_{1X2}\,$ --- сетки по переменным


$X_1$ и $X_2$  с шагом $h_{X1}=\rho_1 N^{-1}_{X1}$ и


$h_{X2}=l_0 N^{-1}_{X2}$ соответственно. Для задачи (\ref{7.3})


строим разностную схему


\begin{multline} \label{7.5}


\La_X Z(X)\equiv \bigg\{\eps^2\bigg\{\, \sum_{s=1,2} A_{ss}(X)


\delta_{\ovl{Xs}\,\widehat{Xs}}+[A_{12}^+(X)(\delta_{X1\,X2}


+\delta_{\ovl{X1}\,\ovl{X2}})+ 


A_{12}^-(X)(\delta_{X1\,\ovl{X2}}+


\delta_{\ovl{X1}\,X2})]+ \\ +\sum_{s=1,2}[


B_s^+(X)\delta_{Xs}+B_s^-(X)\delta_{\ovl{Xs}}]\bigg\} - 


C(X)\bigg\} Z(X) =F(X), \quad X\in D_{1Xh},


\end{multline}


$$


Z(X)=


\begin{cases}


\Phi(X), & X\in \G_{1Xh}\bigcap \G_X;\\


\ovl{Z}_0(X), & X\in \G_{1Xh}\setminus \G_X,


\end{cases}


$$


где $\ovl{Z}_0(X)=\{\ovl{z}_0(x)\}_X$, 


$z_0(x)$, $x\in\Dbh$, --- решение задачи (\ref{3.2}), (\ref{3.4}).





Возвращаясь к переменным $x=(x_1,x_2)$, получаем разностную схему


\begin{equation}


\label{7.6}


\begin{split}


\La z(x)&=f(x), \quad x\in D_{1XhX^{-1}},\\


z(x)&=


\begin{cases}


\vf(x), & x\in \G_{1XhX^{-1}}\bigcap \G;\\


\ovl{z}_0(x), & x\in \G_{1XhX^{-1}\,}\setminus \G.


\end{cases}


\end{split}


\end{equation}


Заметим, что сетка


\begin{equation} \label{7.7}


\Db_{1XhX^{-1}}


\end{equation}


--- сетка на $\Db_1$, вообще говоря, не является ни прямоугольной,


ни равномерной.





Введем функцию


\begin{equation} \label{7.8}


z^*(x)=\left\{


\begin{array}{ll}


z_{(\ref{7.6})}(x), & x\in \Db_{1XhX^{-1}},\\[2mm]


z_0(x), & x\in \Db_{h(\ref{3.4})}\setminus \Db_1


\end{array} \right\}, \quad x\in \Dbh^{\,*},


\end{equation}


где


$$


\Db{}_h^{*}= \Db_{1XhX^{-1}} \bigcup \big\{


\Db_{h(\ref{3.4})}\setminus \Db_1 \big\},


$$


которую назовем решением разностной схемы


(\ref{3.2}), (\ref{3.4}), (\ref{7.6}), (\ref{7.7}).





Заметим, что разностная схема (\ref{3.2}), (\ref{3.4}),


(\ref{7.6}), (\ref{7.7}), аппроксимирующая краевую задачу


(\ref{2.2}), (\ref{2.1}), является безытерационной. При ее


построении используются локально  в области погранслоя


сгущающиеся сетки; шаблон схемы (сетка) подстраивается к


границе области.


\medskip





2. Рассмотрим схему (\ref{3.2}), (\ref{3.4}), (\ref{7.6}), (\ref{7.7}).


Для простоты будем предполагать  выполненным условие


\begin{equation} \label{7.9}


a_1(x)=a_2(x)=a(x), \quad x\in\Db, \quad r(x,\G)\le m,


\end{equation}


т.\,е. коэффициенты при вторых производных в уравнении (\ref{2.2a})


в $m$-окрестности границы $\G$ одинаковы.





Заметим, что схема (\ref{7.4}), (\ref{7.5})


(схема (\ref{7.6}), (\ref{7.7})) --- схема на девятиточечном


шаблоне, вообще говоря, не является монотонной. Условие


\begin{equation} \label{7.10}


\left.


\begin{array}{c}


A_{11}(X)(h_{X1})^{-1}- |A_{12}(X)|(h_{X2})^{-1} \\[2mm]


A_{22}(X)(h_{X2})^{-1}- |A_{12}(X)|(h_{X1})^{-1}


\end{array}


\right\}


\ge 0, \quad \ X\in D_{1Xh},


\end{equation}


является достаточным для $\eps$-равномерной монотонности схемы


(\ref{7.4}), (\ref{7.5}). Это условие будет выполняться, например,


в том случае, когда величина $m$ в $\rho_{1(\ref{5.3})}$


($\rho_1$ определяет множество $\Db_{1(\ref{5.4a})}$) достаточно мала.


Далее предполагаем выполненным условие (\ref{7.10}).





Для простоты полагаем


\begin{equation} \label{7.11}


N_{X1}=N_{X2}=N,


\end{equation}


где $N_{Xs}=N_{Xs\,(\ref{7.4})}$. \ С использованием техники мажорантных


функций (см., напр., \cite{1}) устанавливается оценка


\begin{equation} \label{7.12}


|u(x)-z^*(x)|\le MN^{-2}\ln^2 N, \quad \ x\in \Db{}_h^{*}.


\end{equation}


Таким образом, разностная схема


(\ref{3.2}), (\ref{3.4}), (\ref{7.6}), (\ref{7.7}) при условиях


(\ref{7.9})--(\ref{7.11}) сходится $\eps$-равномерно.





\begin{thms} \label{t.7.1}


Пусть выполняется условие теоремы $\ref{t.3.1}$, а также условие


$(\ref{7.9})$. Тогда решение разностной схемы


{\em (\ref{3.2}), (\ref{3.4}), (\ref{7.6}), (\ref{7.7})} при условии


$(\ref{7.10})$ сходится


$($при $N_1, N_2, N_{X1}, N_{X2} \to\infty)$


$\eps$-равномерно. При дополнительном условии $(\ref{7.11})$ для сеточных


решений справедлива оценка $(\ref{7.12})$.


\end{thms}





\begin{thebibliography}{9}





\bibitem{1}


Шишкин Г.И. {\em Сеточные аппроксимации сингулярно возмущенных


эллиптических и параболических уравнений}. --


Екатеринбург: УрО РАН, 1992. -- 233~c.


\bibitem{2}


Miller J.J.H., O'Riordan E., Shishkin G.I.


{\em Fitted numerical methods for singular perturbation problems}. --


Singapore: World Scientific, 1996. -- 166~p.


\bibitem{3}


Roos  H.-G.,  Stynes  M.,  Tobiska L.


{\em Numerical methods for singularly


perturbed differential equations\/${:}$ convection--diffusion and


flow problems}. -- Berlin: Springer-Verlag, 1996. -- 348~p.


\bibitem{4}


Farrell P.A., Hegarty A.F., Miller J.J.H., O'Riordan E.,


Shishkin G.I.


{\em Robust computational techniques for boundary layers}. --


Boca Raton: Chapman \& Hall/CRC Press, 2000. -- 254~p.


\bibitem{5}


Самарский А.А. {\em Теория разностных схем}. -- М.: Наука, 1989.


-- 616~c.





\end{thebibliography}





\vskip 2cm





\post{\it Институт математики и}{\it Поступила}


\post{\it механики Уральского отделения}{30.10.2001}


\post{\it Российской Академии наук}{}





\label{end}


\end{document}








